\scanpage{43}
% L1–6
(i) $\operatorname{Ext}^{n}(k,A)\simeq k$

\deleted{(\unsure{i bis|ii}) $\operatorname{Ext}^{\unsure{n}}(A/\unsure{\mathfrak n|\mathfrak N},A)\simeq\widehat{\unsure{A|A/\mathfrak n}}$}

\deleted{(ii) $\operatorname{Ext}^{n-1}(k,A)\simeq0$}

(iii) $\operatorname{Ext}^{n+1}(k,A)=0$

(i bis) $\operatorname{Ext}^{n}(A/\mathfrak n,A)\simeq\widehat{A/\mathfrak n}$

\margin{Prouver que chacune des hyp. (i) (iii') (i bis) (vi) implique l'\unsure{équidimensionnalité|équidemensionnalité}.}
\deleted{\unsure{la seule|la seconde} $\lg\operatorname{Ext}^{n}(A/\mathfrak n,A)=\lg(A/\mathfrak n)$}

% L7–13
(iv) Le foncteur $N\mapsto\operatorname{Ext}^{n}(N,\unsure{A|M})$ est dualisant

(iv bis) \rule{.30\linewidth}{.4pt} exact ($N$ de long. finie)

(v) \deleted{\unsure{est injectif|est précisément}} $H^{n}(A)$ dualisant

(v bis) $H^n(A)$ injectif

(v ter) $H^n(A)$ comonogène

(vi) $A/\sum x_iA$ est autodual

(vii) $\overline A\simeq\overline\Omega_A$

% L14–17
\inserted{[signif. $N\mapsto\operatorname{Ext}^{n}(N,M)$ est exact]}
\diagram{\scriptsize\begin{tikzcd}[ampersand replacement=\&,column sep=tiny,row sep=small]
 \boxed{\text{(vii)}\Leftrightarrow\text{(v)}\deleted{\unsure{\text{(v)}|\text{(vi)}}}\Leftrightarrow\text{(v bis)}\Leftrightarrow\text{(v ter)}}
 \arrow[d,Rightarrow] \&
 {\wedge\ \text{$A$ équidim.}\ \deleted{\unsure{\text{(v)}|\text{(iv)}}}} \&
 \boxed{\text{(iv bis)}\Leftrightarrow\text{(iv)}}\arrow[l,Rightarrow]\arrow[d,Rightarrow]\\
 \boxed{\text{(\unsure{iii'|iii})}}\& {}\& {\text{(i bis)}\quad\text{(i)}\quad\text{(vi)}}
\end{tikzcd}}
$A$ est C-M. et $\overline\Omega_A=\overline A$

% L18–23
Prouvons que (vii) implique que $A$ est C-M. \inserted{On peut supposer $A$ \uline{complet}} \inserted{(\unsure{Récurrence|Raisonnement|Rien} non trivial si $n=0$)} En effet, $A$ est équidim. puisque le module $\overline\Omega_A$ l'est. Soit $x\in A$ un diviseur de zéro, il n'est pas diviseur de zéro dans $\Omega_A$. Je dis \deleted{donc} que $\overline\Omega_{A/xA}\simeq\Omega_A/x\Omega_A$

% L24–29
(C'est là \deleted{Ext$_{\!A'}^{\unsure{b|r|*}}$} le résultat spécial à la récurrence). En effet,
\[
 \overline\Omega_{A/xA}=\deleted{\unsure{H^{n-1}}(A/xA)\,\unsure{\operatorname{Hom}|D}}\ \text{or }\operatorname{Ext}_{A'}^{\unsure{n-1|n+1|r-n}}(A/x,A')
\]
or
\[
 \deleted{\unsure{H^{n-1}}(A/xA)^{r-n+1}=\unsure{\operatorname{Hom}}}
\]
\[
 \operatorname{Ext}^{*}_{A'}(A/x,A')\Longleftarrow
 \operatorname{Ext}^{i}_{A}\bigl(A/x,\underbrace{\overset{\inserted{r-n}}{\operatorname{Ext}_{A'}^{j}(A,A')}}_{\Omega_A^{(j)}}\bigr)
\]
soit \deleted{d'où}

% L30–34
\uline{Lemme}
\[
 \Omega^{(*)}(A/x)\Longleftarrow\operatorname{Ext}^{i}_{A}(A/x,\Omega_A^{(j)})
\]
\margin{[Bien connu $=$ pour $A/x$ remplacé par un $A$-module quelconque]}
\deleted{ou \unsure{$A/x$|$Ax$}} \deleted{\unsure{est-il}}

d'où
\diagram{\begin{tikzcd}[ampersand replacement=\&,column sep=small]
 0\arrow[r] \& \operatorname{Ext}_{A}^{1}(A/x,\Omega_A)\arrow[r] \&
 \underbrace{\Omega_{A/x}}\arrow[r] \&
 \operatorname{Hom}^{\deleted{\unsure{0|q}}}_{A}(A/x,\Omega_A^{(1)})\arrow[r] \& 0 .
\end{tikzcd}}

% archive sheet number: 42
